% Tecnicas de dependencia sobre el dia-estacion. Predictores: los tres
% puntajes factoriales de analisis_factorial.tex (reescalados a varianza
% unitaria con los momentos de entrenamiento) mas 12 indicadoras. Las cifras
% vienen de docs/analisis_discriminante.pdf con el 2025 completo.

\subsection{Clasificación: regresión logística y análisis discriminante}
\label{sec:clasificacion_tecs}

Los puntajes factoriales de \cref{sec:analisis-factorial} se unen por
estación y fecha con las variables categóricas del día-estación y con la
respuesta \texttt{excede\_anio\_5}. Los puntajes no tienen varianza unitaria
(desviaciones de 1.14, 1.10 y 1.21 en el ajuste), de modo que se reescalan
con la media y la desviación estándar de 2021--2024, nunca con las de 2025;
así cada razón de momios se lee por desviación estándar del factor y las
tres son comparables entre sí. Las categóricas \texttt{zona},
\texttt{tipo\_dia} y \texttt{temporada} entran como 12 indicadoras contra
los niveles de referencia Noreste, día laborable e invierno. La matriz de
diseño tiene 15 predictores; 16,058 filas entran al ajuste y 2,978 a la
validación.

\subsubsection{Regresión logística}\label{sec:reg_log}

El modelo se ajusta únicamente sobre 2021--2024. Las razones de momios se
reportan en la \cref{tab:momios-log}; ningún intervalo de los tres factores
cruza el 1. El forzamiento fotoquímico tiene el efecto mayor: una desviación
estándar adicional multiplica los momios de excedencia por 2.78, contra 1.92
de la combustión primaria. La lectura es química: los precursores son
necesarios pero la velocidad y el volumen de las reacciones que producen
ozono dependen de la temperatura y la radiación, por lo que la condición
fotoquímica pesa más que la disponibilidad de ingredientes. La ventilación
es el único factor protector (0.90): el viento dispersa.

\begin{table}[t]
  \centering
  \caption{Razones de momios del modelo logístico (\texttt{exp} del
    coeficiente) con intervalo del 95\,\% bajo independencia. Los factores se
    interpretan por desviación estándar; las indicadoras, contra su nivel de
    referencia: Noreste (zona), día laborable (\texttt{tipo\_dia}) e
    invierno (\texttt{temporada}). Todos los coeficientes son significativos
    ($p<0.001$) salvo \texttt{tipo\_dia\_sabado} ($p=0.04$) y
    \texttt{tipo\_dia\_festivo} ($p=0.13$).}
  \label{tab:momios-log}
  \small
  \begin{tabular}{lrc}
    \toprule
    Variable & Momios & IC 95\,\% \\
    \midrule
    Combustión primaria     & 1.92 & [1.82, 2.02] \\
    Ventilación             & 0.90 & [0.86, 0.95] \\
    Forzamiento fotoquímico & 2.78 & [2.62, 2.94] \\
    \midrule
    \texttt{zona\_Centro}    & 2.74 & [2.32, 3.22] \\
    \texttt{zona\_Noroeste}  & 3.84 & [3.32, 4.43] \\
    \texttt{zona\_Norte}     & 1.76 & [1.50, 2.06] \\
    \texttt{zona\_Sur}       & 2.90 & [2.44, 3.44] \\
    \texttt{zona\_Sureste}   & 1.82 & [1.59, 2.09] \\
    \texttt{zona\_Suroeste}  & 3.48 & [3.02, 4.01] \\
    \midrule
    \texttt{tipo\_dia\_sabado}  & 1.13 & [1.01, 1.27] \\
    \texttt{tipo\_dia\_domingo} & 1.43 & [1.27, 1.62] \\
    \texttt{tipo\_dia\_festivo} & 1.17 & [0.95, 1.44] \\
    \midrule
    \texttt{temporada\_primavera} & 3.54 & [3.08, 4.08] \\
    \texttt{temporada\_verano}    & 1.54 & [1.31, 1.83] \\
    \texttt{temporada\_otono}     & 3.04 & [2.66, 3.49] \\
    \bottomrule
  \end{tabular}
\end{table}

Las zonas con mayores momios son Noroeste (3.84) y Suroeste (3.48): a igual
meteorología y calendario, el poniente de la ciudad casi cuadruplica los
momios del Noreste, el ordenamiento espacial que
\cref{sec:contexto-estaciones} anticipó. La primavera (3.54) y el otoño
(3.04) triplican los momios del invierno, que con menor temperatura y
radiación frena las reacciones fotoquímicas aunque los precursores estén
presentes; el verano queda en 1.54. El calendario laboral apenas influye:
solo el domingo (1.43) tiene un efecto claro.

\paragraph{Errores estándar bajo dependencia.} Las 15 estaciones de un
mismo día comparten el clima y los días consecutivos se parecen, por lo que
los errores estándar bajo independencia son optimistas. Se corrigen con el
estimador sandwich agrupado por día (1,460 conglomerados): se inflan entre
1.5 y 2.5 veces, más en las indicadoras de calendario, constantes dentro de
un día, y nada en las de zona, que varían dentro de él. Los intervalos
corregidos de los tres factores son casi el doble de anchos y siguen sin
cruzar el 1 (\cref{fig:momios-corregido}): ventilación $[0.84, 0.97]$,
forzamiento fotoquímico $[2.48, 3.10]$. El único término que pierde
significancia al 5\,\% es el sábado. Agrupar además por estación no es
confiable con 15 conglomerados, pero deja claro que los siete contrastes de
zona descansan en 15 unidades, no en 16,058 filas.

\begin{figure}[t]
  \centering
  \includegraphics[width=\linewidth]{momios_ic_corregido.pdf}
  \caption{Razones de momios de los factores con intervalos del 95\,\% bajo
    independencia (gris) y agrupando por día (color). Los intervalos
    corregidos son más anchos y ninguno cruza el 1.}
  \label{fig:momios-corregido}
\end{figure}

\paragraph{Evaluación y validación.} La evaluación se realiza sobre 2025,
no visto durante el ajuste, con AUC, sensibilidad y especificidad; la
exactitud se descarta porque la tasa base de excedencia en validación es
43.97\,\%. El modelo alcanza un AUC de 0.850 (\cref{fig:roc-modelos}), por
encima del criterio de logro de 0.75. El umbral se elige por el índice de
Youden sobre el ajuste y resulta en 0.282; a ese umbral la sensibilidad es
0.698 y la especificidad 0.853 (panel izquierdo de la \cref{fig:matrices}).
Un umbral de 0.5 ilustra por qué la exactitud no sirve como criterio: la
especificidad sube a 0.964, pero la sensibilidad cae a 0.326, es decir, el
modelo deja de detectar dos de cada tres excedencias reales.

\subsubsection{Análisis discriminante lineal}\label{sec:lda}

La segunda técnica se ajusta sobre la misma matriz de diseño, la misma
respuesta y la misma partición. El contraste es deliberado: el discriminante
supone normalidad multivariada dentro de cada clase y homocedasticidad entre
clases, supuestos que la logística no requiere.

\paragraph{Verificación de supuestos.} Los puntajes factoriales son
combinaciones lineales de predictores con $|\text{sesgo}|<0.5$ tras Box-Cox.
Sobre los puntajes de entrenamiento separados por clase, los seis
coeficientes de sesgo quedan entre $-0.091$ y 0.489, de modo que la
aproximación a la normalidad se sostiene. La homocedasticidad no: no se
aplica la prueba M de Box, que con más de 15,000 observaciones rechaza
diferencias arbitrariamente pequeñas, sino la comparación directa de las
matrices de covarianza por clase. Los días de excedencia son menos dispersos
que los de no excedencia en los tres factores, con razones de desviación
estándar de 0.894, 0.888 y 0.742, y la varianza generalizada baja de 1.081 a
0.346; la discrepancia se concentra en el eje fotoquímico. Un segundo desvío
afecta a las 12 indicadoras binarias, que no son normales bajo ninguna
transformación. El discriminante lineal tolera ambos desvíos en la práctica,
y esa tolerancia se verifica empíricamente en \cref{sec:comparacion}.

\paragraph{Ajuste y evaluación.} Las probabilidades previas se fijan en las
proporciones observadas en el ajuste (0.706 y 0.294) porque la tasa base es
información del problema. Con dos clases hay una sola función
discriminante, y el orden y el signo de sus coeficientes reproducen los de
la \cref{tab:momios-log}: el forzamiento fotoquímico
(0.795) domina sobre la combustión primaria (0.506), la ventilación va en
contra ($-0.109$), el Noroeste y el Suroeste son las zonas de mayor riesgo y
la primavera (1.054) y el otoño (0.772) pesan más que el verano (0.176).
Los coeficientes indican la dirección de máxima separación entre clases; no
son razones de momios. La evaluación usa la
probabilidad posterior de excedencia como puntaje continuo, para que la ROC
sea comparable con la logística. El discriminante alcanza un AUC de 0.846;
al umbral de Youden (0.249) la sensibilidad es 0.758 y la especificidad
0.788 (panel derecho de la \cref{fig:matrices}).

\subsubsection{Comparación de los dos modelos}\label{sec:comparacion}

Las dos técnicas son prácticamente equivalentes sobre estos datos: los AUC
difieren en tres milésimas (0.850 contra 0.846), los puntajes correlacionan
a 0.996 y las clasificaciones coinciden en 93.8\,\% de los días-estación de
validación. La prueba de DeLong para curvas ROC correlacionadas sí detecta
la diferencia ($Z=3.75$, $p=0.0002$, IC 95\,\% $[0.002, 0.005]$): con 2,927
observaciones tiene potencia para tres milésimas, una diferencia
estadísticamente real y sin consecuencia práctica. Donde sí difieren es en
el reparto de errores al umbral de Youden: el discriminante detecta más
excedencias (sensibilidad 0.758 contra 0.698) a cambio de más falsas alarmas
(especificidad 0.788 contra 0.853).

\begin{figure}[t]
  \centering
  \includegraphics[width=\linewidth]{roc_modelos.pdf}
  \caption{Curvas ROC de ambos modelos en la partición de validación (2025).
    Las dos curvas se superponen; el discriminante se traza punteado.}
  \label{fig:roc-modelos}
\end{figure}

\begin{figure}[t]
  \centering
  \includegraphics[width=\linewidth]{matriz_confusion.pdf}
  \caption{Matrices de confusión en validación, cada modelo a su propio umbral
    de Youden. Los porcentajes son proporciones por fila (clase observada).}
  \label{fig:matrices}
\end{figure}

Esta coincidencia es el resultado relevante: la heterocedasticidad
documentada en \cref{sec:lda} es real y aun así no altera el desempeño, de
modo que la frontera de decisión que ambos trazan es la misma y la
interpretación de las razones de momios no depende del supuesto de
normalidad que el discriminante exige y la logística no. Como contraste
adicional, la logística sobre las puntuaciones del ACP alcanza 0.848 de AUC
(DeLong $p=0.44$ contra 0.850): los factores no pierden capacidad
discriminante y ganan interpretabilidad.

\paragraph{Corrección de la tasa de error.} El error aparente por
resustitución (0.229 en la logística, 0.230 en el discriminante) apenas se
mueve con validación cruzada de 10 pliegues por días completos (0.231 y
0.232) ni con el método de Lachenbruch (0.231): con 16 parámetros sobre
15,704 observaciones el optimismo es del orden de $p/n$ y ninguno de los dos
clasificadores sobreajusta. Lo que sí revela la validación cruzada es el
efecto de año: dentro de 2021--2024 el AUC fuera de muestra es 0.80, no
0.85. Por año, cada uno predicho por un modelo que no lo vio, 2021--2023
quedan en 0.76--0.77, 2024 sube a 0.84 y 2025 a 0.85. El desempeño esperado
en un año cualquiera es 0.80; el 0.85 corresponde a los dos años de ozono
alto.

\subsubsection{Limitaciones}\label{sec:limitaciones-clasificacion}

\textbf{Autocorrelación.} Las 18,631 filas día-estación de ajuste y
validación no son independientes: la $n$ efectiva está más cerca de los
1,826 días distintos. La correlación entre estaciones dentro de un mismo día
promedia 0.53 en la respuesta y llega a 0.85 en el factor fotoquímico; la
autocorrelación de la respuesta a un día de distancia, dentro de cada
estación, es 0.49. Los errores estándar se corrigieron por conglomerado de
día y las conclusiones sobreviven; lo que no puede corregirse bien es la
dimensión de estación. El AUC no se ve afectado, porque se calcula sobre
observaciones no vistas.

\textbf{Ausencia de compuestos orgánicos volátiles.} Los ejes se describen
como consistentes con combustión primaria, ventilación o forzamiento
fotoquímico, nunca como una atribución de régimen NOx--COV. Una razón de
momios de 2.78 dice que los días con más forzamiento fotoquímico tienen
mayores momios de excedencia; no dice cuál precursor limita la formación de
ozono, pregunta que no es contestable con estos datos.

\textbf{Alcance explicativo.} Los predictores son del mismo día que la
respuesta: los modelos describen qué condiciones acompañan a una excedencia
y no anticipan el ozono del día siguiente.

\textbf{Validación en un año atípico.} 2025 tiene 44\,\% de excedencia
contra 29\,\% del ajuste. El AUC amortigua el sesgo de tasa base, pero la
sensibilidad y la especificidad se leen sobre un año con más días propensos
a excedencia que el promedio de la serie.
